
% 来源：EPUB part116.xhtml
\section{从更长远的角度看人}

“我不确定这是我在苹果学到的，但这是基于我在苹果时掌握的数据中学到的。这一点就是，我现在对人的看法更长远。换句话说，当我看到有些事做得不对时，我的第一反应不是去指正，而是说，我们正在打造一支团队，这个团队要在未来的十年做一些伟大的事情，而不仅仅是在明年做。所以，我需要做的是帮助把事情办砸的人，让他能够学到东西，而不是去把问题解决。这样做有时候会痛苦，因为我第一本能还是想去解决问题。但对于人的看法更长远，是我在苹果所做的最大改变。”（乔布斯，1992）

这个很牛。我一直有点这个意思，但没这么清晰。（2022-08-18）

% 来源：EPUB part117.xhtml
\subsection{管理者最重要的品质是正直诚信}

\BookQuestion{网友：什么是管理者最重要的品质？}

最重要的是integrity（正直诚信）。（2010-05-20）

\BookQuestion{网友：好像特别强势的领导管辖下的组织，其创新能力和纠错能力会差一些。是不是也是因为这个原因，所以《基业长青》里第五级领导人往往都不特别强势？}

凡是员工见到领导就战战兢兢的公司，时间长了都会出问题的。因为这种公司员工大多都会慢慢变得没有担当，凡事希望交给上级去决定，效率慢慢会降低。

除非这种公司的产品非常聚焦，大老板确实可以照顾到所有大方面，不然早晚会出问题的。我觉得日本企业这些年在有些行业出问题就是这个原因。

\BookQuestion{网友：我感觉乔布斯就特别强势，下属做得不好，直接一顿臭骂。}

所以他累得不行，早早退出了，还好选了库克！（2018-09-22）

\BookQuestion{“迪士尼前CEO鲍勃·伊格尔解释退休原因，他是在感觉自己过于轻视他人的意见后开始思考卸任CEO的。伊格尔说：渐渐地，我开始越来越不愿听取别人的意见，甚至少了一些宽容，或许是因为我对自己有些过于自信，当你有点成绩的时候，就可能发生这种情况。”（新闻）}

“我开始越来越不愿听取别人的意见”，这可能是我见过的最好的辞任CEO的理由！（2021-12-22）

\BookQuestion{网友：您觉得领导者对待员工的两种方式“爱兵如子”和“慈不带兵”，如何平衡比较好？}

我不太喜欢把公司看成军队的想法。顾名思义，“爱兵如子”就是要时刻想着他们的需求，“慈不带兵”指的是该严格要求时要严格要求。没觉得这两点之间有矛盾，所以不该搞平衡吧？（2019-05-19）

% 来源：EPUB part118.xhtml
\subsection{董事会的作用是不让做不对的事情}

\BookQuestion{网友：有没有好的制度来制衡CEO的行为？}

一般来说都是靠董事会，最重要的是企业文化，体现在如何选人上。一旦选错，一般损失都很大。这些年看到不少公司选错CEO的，最典型的就是惠普，当然也包括雅虎。其实阿里当年1688用错CEO的损失也不小，但相对整个公司而言被忽略不计了。（2018-07-02）

\BookQuestion{网友：腾讯游戏收入已占到公司第一比重。假设您是腾讯董事会成员，您会在游戏拓展方面提些什么建议？}

董事会的作用是不让做不对的事，不应该建议什么。（2010-08-30）

\BookQuestion{网友：您现在是不是不管步步高的具体运营了？}

我大概10年前就不在一线了，所以不了解具体运营的东西。我只管做对的事情或者叫发现错误及时纠正，至于如何把事情做对，我不太在行，大把同事比我强。我在公司里大概可以算个顾问吧，没人问我我就没事。（2010-03-03）

\BookQuestion{网友：您对OPPO或步步高还有没有自己的期望和目标？虽然您说不怎么管，只把把大方向。这个大方向，更倾向于您自身的期望吗？还是凭职业经理人们运作，只要不会有什么问题，您也就认可他们对公司发展上的决定，哪怕和您的设想方向有些出入？}

我当然有期望。但目标是大家定的。我在公司的角色是顾问。理论上讲连大方向都不归我管，但我可以提我的意见。

出入总是有的。我在公司里是个反对派，几乎做什么我都会提反对意见。如果连我的反对意见大家都不怕时，做什么我都会放心一些。

我最怕的就是当老板说什么大家都说“好”。那时公司就危险了。

当然，这里的前提是我认为我们很多同事都在许多方面比我强。如果认为自己是公司里最聪明的人的“老板”是很难认同我这个观点的。（2010-04-02）

从5年10年的角度看，CEO至关重要。从10到50年的角度看，董事会很重要，因为董事会能找出好的CEO。从更长的角度看企业文化更重要，因为一个好的企业文化可以维持有一个好的董事会。GE就是一个好例子。好的公司之所以是好的公司，必然会有些好的产品。但所谓好的产品的寿命是非常有限的。所谓好的生意模式可以让好的产品的寿命大幅度提高。（2010-04-04）

\BookQuestion{网友：苹果公司在20世纪90年代陷入困境时董事会决定请乔布斯回来，IBM公司的董事会在IBM出现问题时也一致决定聘请郭士纳为CEO。在这种关键时刻，公司董事会的选择是否也是公司文化的一种体现？}

是的，没有好的企业文化就很难有好的董事会；没有好的董事会，公司将来的CEO人选出问题的概率就会大很多，因为董事会会变得短视。（2013-04-16）

\BookQuestion{网友：如果持有的企业现在不错，但原来的CEO退休了，换了新的最高决策层，您是否会卖出谨慎地观察一下新的决策层？或者认为企业文化跟新的CEO关系不大呢？}

企业文化可以看作是护城河的一部分。换CEO对任何公司都是会有影响的，但企业文化强大的公司往往能选出好的CEO，从而能让公司继续好下去。（2011-01-09）

% 来源：EPUB part119.xhtml
\subsection{最重要的是建立企业文化}

\BookQuestion{网友：作为企业创始人，工作重心要放在何处？}

取决于缺什么。

\BookQuestion{网友：最缺的是什么？}

这个问题问得好，我不知道怎么回答。一般人会问企业什么最重要，我一般回答缺什么什么重要。接下来问缺什么，我还真有点晕啊。无论如何，缺什么的问题应该是经营企业的人自己应该知道的。（2010-10-20）

在公司里显得最重要的那个部门往往是做得不够好的部门。（2017-06-22）

\BookQuestion{网友：“我过去常常认为一位出色的人才能顶两名平庸的员工，现在我认为能顶50名。我大约把四分之一的时间用于招募人才。”乔布斯的这句话对吗？您觉得CEO每天安排时间应该侧重于哪些方面？}

乔布斯是对的。我当CEO的时候并没有一个时间分配表，但大致知道什么是重要的。简而言之，不要让重要的事情变成紧急的事情。（2019-03-16）

\BookQuestion{网友：有什么好方法？}

就是重要的事情要尽早开始做，不要让其变成紧急的事情。（2020-12-07）

\BookQuestion{网友：请问经营公司有哪些需要注意的关键点？具体来说，对于300人规模公司的一把手，需要在哪几件核心的事上持续花时间？有哪几件事，虽然很重要，却应该尽量分给副手？}

我觉得企业领导最重要的是建立企业文化，确立什么东西不能做。如何把事情做对当然也很重要，但只要是在坚持做对的事情，时间就会站在你这边。（2020-10-26）

\BookQuestion{网友：我的企业属于100\%控股，我每天时间都被各类管理的事占满了。我想在未来省下时间来做投资，请问如何既能控制公司发展又不用事无巨细都要我拍板？}

关于授权的问题，能比我理解得更好或做得更好的人确实不多。授权的过程一般是：指示、指导、协商、授权、放权，最后一条是“Neveroutofcontrol”（永远不要失控，只有好的文化才有可能做到这一点）。这里最重要的还是做对的事情，把事情做对。

在能够把握做对的事情的前提下，容忍把事情做对的过程当中所犯的错误。这一点是大多数人做不到而且是很难学会的，就像投资一样，授权也是简单但很难的事情。（2011-01-04）

\BookQuestion{网友：老段化繁为简、放权的能力真的是很少很少有人能匹敌。}

其实是真的有这种意愿的人很少，和能力关系没那么大。（2024-07-10）

% 来源：EPUB part120.xhtml
\subsection{狼性文化最终会输给人性文化}

\BookQuestion{网友：您如何看待狼性文化？张朝阳最近也在反省之前搜狐的好人文化。}

狼性文化最终会输给人性文化。

此“好人”非彼“好人”，你这么看说明你也还没有是非观。（2017-02-26）

我们不用好人坏人这种词去形容做对的事情和把事情做对，容易引起误解。（2017-02-27）

\BookQuestion{网友：您认为文化差异是东西方企业经营理念不同的原因吗？中国不少企业都在加班加点，有些急功近利。}

好企业在哪儿的经营理念都很相近。美国急功近利的公司也很多，中国公司也不都是急功近利。我本人很多年前就很反对没事加班加点的。我一直认为老是强迫加班加点的部门负责人的管理水平有问题，老是强迫加班加点的公司老板的管理水平有问题。（2010-04-14）

我一直都不鼓励加班，但不是说我反对加班。需要时加点班无可非议，但过分加班会影响员工的生活，长此以往也会让竞争力下降的。我总是希望大家能够高效率工作，同时还能好好生活，能多陪点家人孩子啥的。我们接受“慢慢地富起来”的观点，事实上我们做得也还不错啦。（2019-03-22）

\BookQuestion{网友：乔布斯性格不好，心是善的。得知癌症，他的唯一诉求是要活到参加儿子毕业典礼那一天。}

我认识好几个当年在乔布斯身边工作的人，也认识几个跟他合作过的人。没见谁说过不喜欢乔布斯的任何话，除了说他就事论事的时候是非常严厉的。（2024-10-12）

\BookQuestion{网友：《马斯克传》那本书里，他身边的几乎所有工作人员都说他的缺点，说跟他无法相处……我也认识几个曾经在他身边工作过的人，私底下大家会非常含蓄地说点啥。（2024-10-14）}

芒格对马斯克的评价非常中肯，他说马斯克是个被证明了（certificated）的天才，他的智商可能有190，但他自己以为自己是250。现在看来他干的事情有时候真是250啊。马斯克的天赋有时候可以让他做一个不错的报时人，但他太享受报时了。（2018-09-28）

% 来源：EPUB part121.xhtml
\subsection{选人合适性比合格性重要}

\BookQuestion{网友：能不能保证由合适的人来管理公司，会不会出现由不合适的人来管理公司的风险？或者一旦出现由不合适的人管理公司的风险是否可控？}

没有绝对的办法来保证，但如果选人时先看合适性（价值观匹配）会比只看合格性（做事情的能力）要好得多，选中合适的人的概率要大得多。许多公司挑人时，首先看合格性的，概率上就容易出问题。许多公司本身价值观就有问题，选人自然就困难了。不合适的人管理公司的风险很多时候是难以控制的，结果公司很可能就会垮掉。企业文化比较好的公司往往会更快发现问题，所以纠正的机会会大很多，存活下来的概率就会大很多。（2018-06-30）

选人本分诚信比聪明重要，合适性比合格性重要。（2019-09-10）

人才的合适性往往比合格性重要。合适性就是人才对企业文化的认同，合格性就是人才做具体事的能力。合格性是可以通过培训提高的，合适性很难改变。企业文化如何传递的问题建议看韦尔奇的《赢》。（2011-03-26）

\BookQuestion{网友：如何对文化价值观进行有效考核？}

主要靠时间吧。三观不合在一起时间长了一定会知道的。（2020-10-11）

\BookQuestion{网友：如果本分上有瑕疵就不用的话，很多企业和单位可能会面临无人可用的局面。这就是我纠结的地方。}

那只能说明你企业做得不太好，早晚都是麻烦，所以纠结。（2024-06-24）

\BookQuestion{网友：怕犯错的同志，不是好同志？}

做对的事情过程中犯的错和做错的事情导致的错误结果是完全不同性质的。（2022-08-20）

\BookQuestion{网友：您的公司有没有请空降兵？}

专业人士很多，CEO位置没有。（2010-08-25）

\BookQuestion{网友：大道分享一些看人小技巧或者方法？比如库克决定聘请那位之前担任巴宝莉CEO的女士最后证明是错误的。而大道一开始就认为库克选错了人。}

我没说库克选错人吧，我只是觉得从苹果之外选人，我有点怀疑企业文化的匹配性。事实证明我的怀疑是有道理的。（2023-07-22）

\BookQuestion{网友：有人说惠普的董事会之所以会选择菲奥莉娜，是被其面试时的表演能力所迷惑。}

找空降兵其实不是上策。惠普要走这一步实在可惜。（2010-08-24）

韩非子曰：“宰相必起于州部，猛将必发于卒伍。”雅虎这么老从外面找CEO的办法，说明美国雅虎董事会确实很烂啊。（2011-09-08）

\BookQuestion{网友：工作上长时间与OPPO和vivo的代理公司有合作，这些代理公司也都沿袭了OPPO和vivo的企业文化，本分的行事风格在能接触到的所有人身上都有很好的体现，在手机品牌百花齐放、竞争最激烈的10多年前，那时的步步高和OPPO代理公司是显得有些“格格不入”的一股清流。业务人员基本都是高学历人员，新人基本都是校招的，有一股书生气。几乎没看到他们用其他公司跳槽过来的人（当时业务人员在各个公司跳槽很普遍），其他公司业务人员大多有点江湖气。OPPO和vivo的人员沟通都很坦诚，诚实也务实，各种细节做得相当细，不说大话假话，价格明码实价……这也许就是大道说的本分，今天看这些好像都是应该的，但当时其他品牌为了出货量讲虚话假话的情况是很普遍的。}

校招新人是我们最早的策略，因为觉得企业文化很重要，从其他企业招过来的人往往带着其他企业的文化，认同不容易。前些年为了发展快点我们公司好像确实从其他企业招了不少人，好像也吃了些苦头。欲速不达其实就是这个意思。（2024-08-27）

\BookQuestion{网友：该怎么安排资历老但是能力有限的早期员工呢？}

马云安排得好，叫心软刀子快。（2011-11-03）

\BookQuestion{网友：我觉得裁员对于管理层来讲是个很艰难的事情，但有时候又不得不做，怎么处理这个问题？}

裁员也许是每天都该做的事情，或者说每年都要做的事情。不合适的人早点让他们离开对大家都好，不要等到不得不做的时候才做。裁员大概是最难的一件事情，尽量不要因为短期的困难去裁员。

\BookQuestion{网友：如果个别同事不符合我们的文化，我们也找不到好办法改变他，该怎么处理？}

不管怎么样，如果融不进这个文化，你不让他走他也会想走的。（2019-03-18）

\BookQuestion{网友：从经营企业的角度看，您认同10\%末位淘汰吗？}

个人不太喜欢这种淘汰的办法：如果都是认真招进来的不错的员工，如何才能挑出10\%赶走呢？（2016-04-28）

\BookQuestion{网友：大道怎么看相对考核法（一群人进行排列，前10\%打A，后5\%打C，中间是B）。感觉在相对考核法下跟同事的竞争大于协作，在这种环境下做得不是很开心。步步高的方法是什么呢？}

我当CEO时没这么做过，现在应该也不是这么做的。我记得曾经看过很多类似的东西，但始终觉得有点不踏实，所以没用过。（2020-10-10）

\BookQuestion{网友：您如何看待离职员工在平台或知乎上发表对公司的负面评价？}

这是公司收集意见然后改进的好渠道。有则改之无则加勉吧。（2019-03-16）

% 来源：EPUB part122.xhtml
\subsection{公平心是基本原则}

公平的定义其实非常复杂，有兴趣大家可以看看网易的公开课。

我个人理解，其实没有绝对的公平——就是所有人都觉得公平的公平。

作为游戏规矩的制定者，最重要的是要有公平心。

公平并不等于平等，也不是平均主义。

看某样东西是否公平往往从旁观者的角度会比较客观些。

公平同时也是很难的一件事，需要很多的功夫，比如一个公司的绩效考核体系等。客观讲，我自己没有能力去阐述什么是公平，也许这是大学教授才能干好的事。不过，我可以举一些我们这些年的经历或我们制定的小规矩来说明我们对公平以及平等的理解。

在美国，机场安检有时很严，有些人偶尔会被抽到详细检查（包括搜身）。有时候，我们会看到被抽到的人很不高兴，觉得为什么是他而不是别人，觉得那不公平。我被抽到是感觉很坦然，也很配合，觉得那是公平机制下的结果，也是为了我们每一个人好。

以前经常看到很多公司会有各种规定，比如关于什么人可以分什么样子的房子，什么级别可以配什么样的车，什么级别以上的人出差可以坐头等舱或商务舱等等。我们公司里不分房子、不配车（都是自己买），任何人出差都可以坐头等舱（但公司只报销经济舱）。记得刚开始有这个规矩时，有一次我发现我的头等舱机票被全报了，就问财务为什么。回答是公司只有我可以报头等舱，然后我说，这是不把我当人看，因为我们的规矩是对所有人的。后来就都一样了。

我记得曾经和一家公司的某部门头讨论奖金问题。由于他们部门利润好，不敢按全额把奖金发下去，另外也怕奖金的分配会不公平。我个人认为所谓的奖金其实也是某种契约，公司把一定比例的获利分给员工，封顶实际上是违约的。而由于怕分不好而不发则是最大的不公平。不过，要发得好也确实是一件非常难的事，不光要有很好的绩效考核体系，同时可能也需要很多年经验的积累，但不发肯定是错的。（2011-02-21）

\BookQuestion{网友：出差可以坐头等舱是对的，休息好工作效率就会高，我发现有些公司都不可以坐头等舱，我就觉得这是个问题！}

那你没看懂我说的是什么。我的规定是所有人一视同仁，想坐头等舱（或者商务舱）就自己掏钱。我就一直坐头等舱，也一直自己掏钱，因为我不想为难自己，但同时我也不喜欢等级概念。（2022-02-22）

\BookQuestion{网友：强制而公平。}

这里没有啥强制的概念吧？这里其实强调的是没有等级概念，不是什么级别的就可以坐商务舱或者头等舱啥的，毕竟薪水待遇都已经尽量表达清楚了，其他开销就应该自己来决定。其实按我的概念，酒店标准也该一视同仁，所有人的报销标准都应该一样。

\BookQuestion{网友：我经营企业4年，发现员工对已有的规则和流程不严格遵守，屡屡犯错，可实际规则是由下至上订出来的。现象举例，上班不打考勤机，以忘记为由事后去打考勤。针对上述现象，该怎么应对？}

我觉得大家都不遵守的规矩有可能会是不合理的规矩。小部分人不遵守规矩是执行者的问题。一般而言规矩定得一定要有道理、有意义并且尽量简单和易于执行，比如这个博客博友发言的规矩。

一个组织如果有大家都不执行的规矩存在，长期来讲就是鼓励大家不执行制定的规矩，未来大家的麻烦可就大了——比如国内常见的没有交警和红绿灯的十字路口。面对的方式其实也很简单，就是执行你的规矩或者废掉执行不了的规矩，最好不要有中间地带，不然大家会很忙的。

举个我们自己的小例子。有一年，我回了一趟公司，发现总部大楼前那些没有停车线的地方也都停满了车。我就问我们管后勤的同事，为什么这么多车停在“不可以”停车的地方？他告诉我说，因为“没地方”停。后来他们把那些原来“不可以”停车的地方都画上了停车线，这个违规的问题就解决了。（2011-01-05）

\BookQuestion{网友：你们公司违反了就要开除的规矩有哪些呢？}

涉及开除的条例应该是不多的。我印象中有过因为贪污被抓起来的。（2020-12-03）

\BookQuestion{网友：您对采购员拿回扣是什么看法？}

采购员拿回扣应该开除，给回扣的供应商应该列入禁入对象。（2020-12-04）

\BookQuestion{网友：考核员工短期拉客户数量，必然导致盲目寻找不诚信的人。我认为应该增加一个长效指标：培养的客户产生了多少销售额？引导业务人员更注重客户的质量。}

呵呵，指标的设定实际上要复杂很多。比如如果你用客户的销售额做指标的话，可能就会有人和客户联手做假销售额。我记得以前有时候为了找一个好指标要花几年时间。（2011-03-06）

\BookQuestion{网友：遇到一个工作难题，公司规定配套厂家未按计划完成，每缺一件罚200元，造成生产缺件每台500元，每天有100多台未完成计划，这样持续罚款，厂家挣的钱不够罚款，公司领导也就不罚了。如果制度没有持续性，对配套商也就没有约束了，段总有何建议？}

是自己公司内部的配套厂家还是外部的？如果是外部的，也许你们要有个备用的厂家，内部的也一样。如果定的规矩不实施，比没有还糟，但如果因为定得不合理而实施不了的规矩就一定要改规矩，不然公司会很麻烦。（2011-05-03）

\BookQuestion{网友：我是一个30多岁的研发部门基层管理者，自己技术能力尚可，但在管理下属方面，有时候会对下属过于苛刻，似乎在一些无关紧要的细节上非常执着（事后来看），要求过高搞得下属反感；但有时候又感觉自己很没有原则，甚至在一些重要的事情上容易妥协。}

原则问题上确实不应该妥协，其他时候需要将心比心。我有时候也会很严厉，但都是就事论事。另外，我觉得自己做不到的事情我也很难要求别人（我对小孩也这样，我小时候做不到的事情我是不要求他们的）。不过，管理公司还是要有规矩的，这就需要你找到办法。比如我非常不喜欢一早去上班，所以早年我当CEO的时候，也不好意思要求大家早上班，于是就让各个部门自己定一个上班时间，前提是所有部门至少要有4个小时以上的交叉时间，结果除了我以外大家都选择了8点上班，但我却可以心安理得地10点以后才出现。（2024-10-09）

\BookQuestion{网友：公司最早的一些政策文档应该都是您还在公司时制定的，比如每年利润股权分配、代理商打分考核。最初第一版本设定时，您是如何考虑可以让规则的基础那么简单清晰，同时兼顾效率和公平性的？}

简单清晰才有效率啊，公平心是基本原则。这里公平心非常重要！公平心和公平性是有区别的。最高领导制定规矩一定要花够时间，一定要想长远，要防止大漏洞。修改时一定要有提前量，承诺的东西一定要执行，所以绝对不能轻诺……很久没在一线，一下想不起来还有啥要注意的，但基本原则依然是想长远！（2022-02-24）

\BookQuestion{网友：人有恒产，就有恒心，然后生活就踏实。}

其实企业也一样，只要企业内的规矩是靠谱的，员工就会有恒心，我当年就不会离开小霸王，也许阿里这个生意就是我们在干着，马云可能就要去干别的了（玩笑哈）。（2013-03-03）

\BookQuestion{网友：在不靠谱的公司待过才知道啥应该是靠谱，跟不靠谱的人交往过，才知道啥是靠谱的人。}

其实只有在靠谱的地方待过才知道什么叫不靠谱，反过来其实不成立。（2013-03-06）

\BookQuestion{网友：创业公司中如果有配偶或亲戚或朋友在一起工作，您觉得长期看有哪些好处和坏处？}

创业公司里如果亲戚朋友一堆，可能短期对凑齐人有好处，长期坏处比较多，会形成裙带关系，而且不容易留住别的人才。（2020-12-05）

我觉得在我们中国文化的大前提下，如果允许直系亲属在公司的话，那10年8年后重要岗位基本上就都是那谁谁谁的亲属了，那企业怎么可能还有竞争力？再说，有了这条规矩，我们同事们可是省了无穷多的麻烦啊，不然每个稍有点权力的人都会背上很多人情债的。

\BookQuestion{网友：您认为一个企业的采购等敏感部门都是大股东家属，是否会有些问题，这样的企业能投资吗？}

从好的地方讲那可能是大股东找不到他信任的别人。从不好的地方讲，他有可能通过这个为自己牟利。我个人认为从长期而言，还是尽量不要用亲属的好。很难单用这个作为是否投资的标准。（2011-02-21）

% 来源：EPUB part123.xhtml
\subsection{钱是保健因子，不是激励因子}

\BookQuestion{网友：管理人才有什么重要原则，可以分享一下吗？}

据说马云说过，人待不住就两个原因，钱给少了或者是心委屈了。马云这句话非常经典，把保健因子和激励因子说得非常形象。钱是保健因子，不是激励因子，给多了没用，给少了不行（绝大多数人在这点上都错了，因为钱不是万能的，当然没钱也是万万不能的）。激励因子实际上就是不要让人心委屈了，要让人觉得有意思。（2019-10-19）

经常看到报道说有人发多少多少钱去激励员工，甚至据说还有扛着钱到球场边督战的。

不知道是不是真的重赏之下必有勇夫，但在企业经营中靠多发钱对员工其实是没有激励作用的。

钱其实是保健因子，多了没用，少了不行的东西。比如，如果员工已经很努力在工作了，多发钱并不会让他们更努力，但少发是会留不住人的。所以，发多少其实体现的是公平。

公司应该按承诺发给员工应该得的报酬，多发少发其实都会破坏公平。

某些所谓老板给员工发红包的说法其实也是有点问题的，因为其实这个红包就是员工应该得的。记得早期还在小霸王时，由于公司人还很少，很多同事的年终奖都是我自己一个一个发到大家手里的。

有时有人会说“谢谢老板”，我就会很严肃地说应该是公司谢谢大家一年的努力。这个真不是客气。（2011-02-17）

\BookQuestion{“传OPPO员工离职后照发12万年终奖，主动离职、辞退都N + 1。”（新闻）}

过去20多年都是如此，那是契约精神里的东西，没什么好说的。（2019-04-20）

\BookQuestion{网友：您对股权激励计划怎么看？}

实际上是保健计划，可人们非要叫激励计划。（2013-06-26）

你如果不理解股票激励政策的好处，只要想想没有股票激励的坏处就明白了。如果员工的工作和结果无关的话，最后的结果就一定是老式国营企业的样子。（2011-05-28）

\BookQuestion{网友：在小公司里给核心员工分股权有没有实际意义？}

你是说谷歌的初期还是苹果的早期？（2015-02-25）

\BookQuestion{网友：当年步步高给员工分配股权的时候，是用现金或者股息来购买股份的。能否介绍一下具体怎么分的？}

我离开一线20多年，很多细节已经不太记得了。应该没有所谓“分配股权”的时候，但早期确实有过允许一定级别以上的员工用年终奖的一部分（好像是20\%）自愿购买公司的股份。（2022-07-05）

这里给个小建议：股权激励（实际上是个保健因子）需要高管们拿钱买股票的时候，可以用期权的形式，或者公司借钱给他们（用未来的收入分期还）。（2024-10-29）

我们30年前就这么干了。（2024-10-30）

\BookQuestion{网友：怎么对待离职员工的股权，如有员工干了几年离职，企业付费把股权回购，违背本分吗？}

如果是有言在先没问题，这是双方的契约精神。（2019-03-22）

\BookQuestion{网友：国内很多公司给员工派发的股票越来越少，搞得新人很心寒，这样没有新人的动力，公司能持续发展吗？}

呵呵，如果你反过来想想，比如给员工发的股票越来越多会怎么样？这个问题看看那些百年老店怎么做的就明白了。

《从优秀到卓越》里好像讲到过这个问题。好公司的收入一般都不是最高的。我在中欧EMBA时学过一个很重要的东西，就是钱是保健因子而不是激励因子。（2010-03-28）

\BookQuestion{网友：请问在分配奖金时，需要遵循什么标准才能最大限度调动员工的积极性？对于业绩突出和业绩最差的员工，如何拉开收入差距？是公开还是不公开？}

钱是保健因子，就是少了不行的意思，奖金并不能调动大家的积极性。

奖金实际上应该是种契约，是跟业绩挂钩的（所以不是老板赐给大家的），比如在某些条件前提下，企业获得的利润的20\%给员工。分配时则是根据你的绩效考核体系去分的。

不宜公开。公平指的是公平心，但从每个人自己的角度看，未必就觉得公平。不过，如果绩效考核体系完整的话，其实也和公开没太大差别了。（2019-04-26）

\BookQuestion{网友：经营企业赚到了钱，如何分钱给合伙人及员工，应该基于什么去分配？}

应该按照事先说好的原则吧？最重要的是要有公平心。（2020-11-30）

\BookQuestion{网友：您对挖竞争对手的人才怎么看？}

我们公司一般不主动挖人，但招人总是可以的。

一般整编制挖人的短期效率高，但长期对企业文化传承不利，而且经常会出现损人利己的事情，最后倒过来破坏整个行业的生态环境，最后也会伤害到自己公司。自己培养人短期看起来慢，但长期看往往是效率最高的办法。不过，老是被挖人的，尤其是关键人物老被挖的企业，一定要好好检讨自己的保健因子是不是不够，为什么会留不住人呢？由于机制的问题，国有企业要留住人才确实困难大很多，好的人才不被人挖也有很高的机会自己流走。（2011-03-08）

\BookQuestion{网友：您在步步高的时候，有没有定期找时间亲自培训下面的人？您怎么看待员工培训？}

培训的成本很贵，但不培训的成本更贵。不知道第二点的企业走不远。（2013-10-28）
